\documentclass[11pt]{article}
\usepackage{amsmath,amssymb,amsthm}
\usepackage[T1]{fontenc}
\usepackage[utf8]{inputenc}
\usepackage{amsmath,amssymb,amsthm}
\usepackage[margin=1in]{geometry}
\usepackage{microtype}
\usepackage[hidelinks]{hyperref}

\newcommand{\Z}{\mathbb{Z}}
\newcommand{\F}{\mathbb{F}_2}
\newcommand{\surg}[2]{S^3_{#1}(#2)}
\DeclareMathOperator{\rk}{rk}
\DeclareMathOperator{\intt}{int}

\newtheorem{theorem}{Theorem}
\newtheorem{lemma}[theorem]{Lemma}

\newcommand{\fptitle}{Review of Problem 6 -- Astra}
\newcommand{\reviewtag}{Reviewer 3}
\usepackage{fp-review}

\begin{document}
\maketitle

\section{Review}

\paragraph{Correctness.}


\begin{itemize}
	\item  1: essentially flawless mathematical argument
	\item  0.5: minor errors that are easily corrected
	\item  0: blank solution, or acknowledged partial result
	\item  $-0.5$: major errors
	\item  $-1$: proposed proof makes no progress towards the result.
\end{itemize}
Note that points should be docked for wrong statements/ solutions.

\noindent Score: 1.

\paragraph{Exposition.}


\begin{itemize}
	\item  1: well-written and clear solution.
	\item  0: argument that is possible to follow with serious effort
	\item  $-1$: very poor exposition
\end{itemize}

\noindent Score 1.

\paragraph{Attributions.}
\begin{itemize}
	\item  1: solutions contains fully appropriate attributions
		to related work and theorems being used.
	\item  0: moderate issues with missing attributions or hallucinated references
	\item  $-1$: blatant plagiarism or hallucinations
\end{itemize}

\noindent Score: 0. 

 I didn't find any errors in the math, though it's convoluted at times.

The argument can be followed without "serious effort", hence a score of 1. That said, the exposition is quite poor. Sometimes it contains weird and unnecessary phrases. Other times it provides wordy or confusing explanations (e.g. ``by the surgery homology calculation and the
connected-sum homology formula'') for somewhat straightforward facts. 

The attributions are mostly fine, but there are repeated instances of citing secondary sources. 

\newpage
\section{Submission with inline comments}


\begin{abstract}
Every Dehn surgery on an \(L\)-space knot in \(S^3\) has at most two
nontrivial prime summands. We prove this using established results about
reducing slopes and prime decompositions, together with a
group-theoretic genus bound for integer homology sphere summands of
surgeries on fibered knots.
\end{abstract}
\review{Abstract is poorly phrased. For instance, it sounds like it's referring to "genus bounds for homology spheres" when presumably it means genus bounds for the knots}

\section{Statement and conventions}

We use the standard ambient-manifold convention: the knots under
consideration lie in \(S^3\). 
\review{Awful opening sentence, knots living in $S^3$ is hardly a convention.}

A rational homology sphere \(Y\) is an
\(L\)-space if
\[
    \dim_{\F}\widehat{HF}(Y;\F)=|H_1(Y;\Z)|.
\]
An \(L\)-space knot is a knot admitting a nonmeridional \(L\)-space
surgery. We allow either sign for this surgery coefficient; when the
coefficient can be chosen positive, we explicitly say that the knot
admits a positive \(L\)-space surgery.

For a knot \(K\), write \(\surg{p/q}{K}\) for Dehn filling along
\(p\mu+q\lambda\), where \(\mu\) is a meridian and \(\lambda\) is the
Seifert longitude. For finite slopes we take \(\gcd(p,q)=1\) and \(q>0\);
the meridional slope is denoted by \(\infty\).

All homology groups without specified coefficients have coefficients
in \(\Z\). Prime summands mean the non-\(S^3\) factors in the
Kneser--Milnor prime decomposition
\cite[Theorem~1]{Milnor}. In particular, \(S^3\) contributes no
nontrivial prime summands.

\begin{theorem}\label{thm:main}
Let \(K\subset S^3\) be an \(L\)-space knot. For every
\(r\in\mathbb{Q}\cup\{\infty\}\), the manifold \(\surg{r}{K}\) has at most
two nontrivial prime summands.
\end{theorem}

We first state precisely the published surgery results that will be
used, and then prove the required group-theoretic estimate.

\section{Inputs from surgery theory}

\begin{theorem}[Fiberedness]\label{thm:fibered}
If \(K\subset S^3\) admits a positive \(L\)-space surgery, then its
exterior fibers over \(S^1\), with fiber a Seifert surface of genus
\(g(K)\).
\end{theorem}

This is the fiberedness assertion of
\cite[Theorem~1]{BaldwinSivek}, which supplies a proof of the result.

\review{This is weirdly phrased, it mixes the content of the BS theorem with the definition of a fibered knot}

\begin{theorem}[Hom--Lidman--Zufelt]\label{thm:slope}
Suppose that a nontrivial knot \(K\subset S^3\) admits a positive
\(L\)-space surgery and is not a cable knot, where nontrivial torus
knots are counted as cables of the unknot. If \(\surg{r}{K}\) is
reducible, then
\[
    r=2g(K)-1.
\]
\end{theorem}

\review{Again, this is correct but weirdly incompetent: ``nontrivial torus
knots are counted as cables of the unknot''... What is ``are counted as'' doing here? Torus knots simply {\it are} cables of the unknot.}

We use the non-cabled formulation stated in the abstract of
\cite{HLZ}. Their Theorem~1.3 establishes the hyperbolic case, and
their Section~1 explains the reduction from the non-cabled case using
the established torus- and satellite-knot results.

\begin{theorem}[Howie]\label{thm:three}
If a manifold obtained by Dehn surgery on a knot in \(S^3\) has more
than two nontrivial prime summands, then it has exactly three.
Its prime decomposition has the form
\[
    L_1\#L_2\# Z,
\]
where \(L_1,L_2\) are nontrivial lens spaces and \(Z\not\cong S^3\) is
an integer homology sphere.
\end{theorem}

This is \cite[Corollary~5.3]{Howie}, including the consequences of the
results of Valdez-S\'anchez and Sayari incorporated in that corollary.

\review{missing citations to the original sources}

We will also use the Poincar\'e theorem: a closed simply connected
\(3\)-manifold is homeomorphic to \(S^3\); see \cite{MorganTian}.

\review{again, maybe Perlman should get a shoutout here? LOL}

Consequently, a closed orientable \(3\)-manifold with cyclic
fundamental group has at most one nontrivial prime summand.
Indeed, a decomposition into two non-\(S^3\) factors would, by
Seifert--van Kampen and the Poincar\'e theorem, express its fundamental
group as a free product of two nontrivial groups. Such a free product
is not cyclic: nonidentity elements from its two factors do not
commute, by the reduced-word normal form.

\begin{lemma}[The cable case]\label{lem:cable}
Every Dehn surgery on a nontrivial cable knot in \(S^3\), including a
nontrivial torus knot, has at most two nontrivial prime summands.
\end{lemma}

\review{A poor stylistic choice to single out torus knots.}

\begin{proof}
It suffices to consider a nonprime surgery.
\review{Again, weird opener. Should just start by setting up the notation, as is done in the next paragraph.} 
The classical results for
cable knots imply that such a surgery is along a cabling slope.
For a nontrivial companion, this follows from
\cite[Corollary~4.4 and the proof of Corollary~4.5]{Scharlemann};
the torus-knot case is contained in Moser's classification
\cite{Moser}.

Thus, for some knot \(J\), we may write
\[
    K=C_{a,b}(J),\qquad r=ab,
\]
where \(a\geq 2\), \(\gcd(a,b)=1\), and \(a\) is the winding number
of the cable. The cable-surgery formula
\cite[Corollary~7.3]{Gordon} gives
\[
    \surg{ab}{C_{a,b}(J)}
       \cong L\#\surg{b/a}{J},
\]
where \(L\) is a lens space with fundamental group of order \(a\).
In particular, \(L\) is prime by the preceding observation.

\review{This is a well-known fact about lens spaces.}

Since \(a\geq 2\) and \(\gcd(a,b)=1\), the slope \(b/a\) is
nonintegral. If \(J\) is nontrivial, the nonintegral-surgery theorem of
Gordon--Luecke \cite{GordonLuecke} says that
\(\surg{b/a}{J}\) is irreducible, hence prime.
If \(J\) is the unknot, its exterior has cyclic fundamental group,
and the fundamental group of every filling is a quotient of that
group. Thus \(\surg{b/a}{J}\) again has at most one nontrivial prime
summand. The displayed connected sum therefore has at most two.
\end{proof}

\section{A genus bound for homology sphere summands}

For a finitely generated group \(G\), let \(\rk(G)\) denote the
minimum size of a generating set. This is group-theoretic generating
rank, not the rank of the abelianization.

\review{The above sentence is pure slop.}

\begin{lemma}\label{lem:genus}
Let \(K\subset S^3\) be a fibered knot of genus \(g\), and let
\(p/q\) be a reduced finite slope with \(p\neq 0\).
Suppose
\[
    \surg{p/q}{K}\cong X\# Z,
\]
where \(Z\not\cong S^3\) is an integer homology sphere.
Then
\[
    |p|\leq g.
\]
\end{lemma}

\begin{proof}
Set
\[
    E=S^3\setminus\intt\nu(K),\qquad
    Y=\surg{p/q}{K},\qquad
    G=\pi_1(E),\qquad
    \Gamma=\pi_1(Y),\qquad
    n=|p|.
\]
\review{Or, you could just use $|p|$...}

For a group \(A\), write \(A'=[A,A]\).

\medskip
\noindent\emph{First, the commutator subgroup of \(\Gamma\) has at most
\(2g\) generators.}
The fibration of \(E\) has fiber a once-punctured genus-\(g\) surface
\(F\). Its homotopy exact sequence gives
\[
    1\longrightarrow\pi_1(F)
      \longrightarrow G
      \longrightarrow\Z
      \longrightarrow 1.
\]
The homomorphism to \(\Z\) is the abelianization: a meridian maps to a
generator, and \(H_1(E)\cong\Z\). Hence
\[
    G'\cong\pi_1(F)\cong F_{2g},
\]
where \(F_{2g}\) is the free group of rank \(2g\).

Seifert--van Kampen identifies the filling homomorphism as the
surjection
\[
    G\twoheadrightarrow
    G/\langle\!\langle\mu^p\lambda^q\rangle\!\rangle
    \cong\Gamma.
\]
A surjective homomorphism maps the commutator subgroup onto the
commutator subgroup: every commutator in the target is the image of
a commutator of lifts.
\review{verbose and unnecessary}
Therefore \(G'\twoheadrightarrow\Gamma'\),
and
\begin{equation}\label{eq:upper-rank}
    \rk(\Gamma')\leq 2g.
\end{equation}

\medskip
\noindent\emph{Next, the homology sphere summand forces at least \(2n\)
generators.}
The standard surgery homology calculation gives
\begin{equation}\label{eq:surgery-homology}
    H_1(Y)\cong\Z/n\Z.
\end{equation}
Indeed, \(H_1(E)\) is generated by the meridian, the longitude is
null-homologous, and filling imposes the relation \(p[\mu]=0\).

Because \(Y=X\# Z\) and \(H_1(Z)=0\), the connected-sum homology
formula gives
\[
    H_1(X)\cong\Z/n\Z.
\]
Let \(\widetilde X\to X\) be the connected \(n\)-sheeted covering
corresponding to
\[
    \ker\bigl(\pi_1(X)\longrightarrow H_1(X)\bigr).
\]
For \(n=1\), this is simply the identity covering.

Choose the balls used to form \(X\# Z\), and write
\[
    X^\circ=X\setminus\intt B^3,\qquad
    Z^\circ=Z\setminus\intt B^3.
\]
The inverse image of the ball in \(X\) consists of \(n\) disjoint
balls in \(\widetilde X\). Attach one copy of \(Z^\circ\) to each of
the resulting boundary spheres. This constructs a connected
\(n\)-sheeted covering
\[
    \widetilde Y\longrightarrow Y
\]
with
\begin{equation}\label{eq:cover-sum}
    \widetilde Y\cong
    \widetilde X\#
    \underbrace{Z\#\cdots\# Z}_{n\text{ copies}}.
\end{equation}
To identify the covering subgroup, observe that
\[
    \pi_1(Y)\cong\pi_1(X)*\pi_1(Z).
\]
The covering monodromy is abelianization on the first factor and is
trivial on the second, since \(H_1(Z)=0\). 
\review{Not well justified}
It is therefore precisely
the abelianization
\[
    \Gamma\longrightarrow H_1(Y)\cong\Z/n\Z.
\]
Consequently,
\[
    \pi_1(\widetilde Y)\cong\Gamma'.
\]

The group \(B=\pi_1(Z)\) is perfect because \(H_1(Z)=0\).
It is nontrivial by the Poincar\'e theorem and the assumption
\(Z\not\cong S^3\). A nontrivial perfect group cannot be generated by
one element, since a one-generator group is abelian. Thus
\[
    \rk(B)\geq 2.
\]

Applying Seifert--van Kampen to \eqref{eq:cover-sum} gives
\[
    \Gamma'\cong
    \pi_1(\widetilde X)*
    \underbrace{B*\cdots*B}_{n\text{ factors}}.
\]
Grushko's rank theorem states that generating rank is additive under
free products of finitely generated groups; see \cite{Stallings}.
\review{Grushko should get a citation}
It follows that
\begin{equation}\label{eq:lower-rank}
    \rk(\Gamma')
      =\rk\bigl(\pi_1(\widetilde X)\bigr)+n\rk(B)
      \geq 2n.
\end{equation}
Combining \eqref{eq:upper-rank} and \eqref{eq:lower-rank} yields
\(2n\leq 2g\), and hence \(n=|p|\leq g\).
\end{proof}

\section{Proof of the main theorem}

\begin{proof}[Proof of Theorem~\ref{thm:main}]
Mirroring changes surgery coefficients by a sign and reverses the
orientation of the filled manifold:
\[
    \surg{-r}{\overline K}\cong-\surg{r}{K}.
\]
The \(L\)-space condition is invariant under orientation reversal,
as is the number of prime summands. We may therefore assume that
\(K\) admits a positive \(L\)-space surgery; compare
\cite[Section~1]{HLZ}.

Meridional surgery gives \(S^3\), so assume \(r\) is finite.
If \(K\) is the unknot, its exterior has cyclic fundamental group,
and every filling has cyclic fundamental group. The observation
following Theorem~\ref{thm:three} settles this case.
If \(K\) is a nontrivial cable, Lemma~\ref{lem:cable} settles the
claim.

We may consequently assume that \(K\) is nontrivial and non-cabled.
Suppose, for a contradiction, that
\[
    Y=\surg{r}{K}
\]
has more than two nontrivial prime summands. In particular, \(Y\) is
reducible. Put \(g=g(K)\). By Theorem~\ref{thm:slope},
\[
    r=n=2g-1.
\]
Since \(K\) is nontrivial, \(g\geq 1\), so \(n>0\).

Theorem~\ref{thm:three} supplies a decomposition
\[
    Y\cong L_1\#L_2\# Z,
\]
where \(L_1,L_2\) are nontrivial lens spaces and \(Z\not\cong S^3\)
is an integer homology sphere. Let
\[
    a_i=|H_1(L_i)|,\qquad i=1,2.
\]
Both \(a_i\) exceed \(1\). The surgery homology calculation and the
connected-sum homology formula imply
\[
    n=|H_1(Y)|=a_1a_2>1.
\]
Together with \(n=2g-1\), this shows that \(g\geq 2\).

By Theorem~\ref{thm:fibered}, \(K\) is fibered.
Apply Lemma~\ref{lem:genus} to the integer homology sphere summand
\(Z\), taking \(X=L_1\#L_2\). It gives
\[
    2g-1=n\leq g.
\]
This is impossible for \(g\geq 2\). The supposed decomposition cannot
exist, proving the theorem.
\end{proof}

\begin{thebibliography}{99}

\bibitem[BS22]{BaldwinSivek}
J.~A. Baldwin and S. Sivek,
\emph{\(L\)-space knots are fibered and strongly quasipositive},
The Open Book Series \textbf{5} (2022), no.~1, 81--94.

\bibitem[Gor83]{Gordon}
C.~McA. Gordon,
\emph{Dehn surgery and satellite knots},
Trans. Amer. Math. Soc. \textbf{275} (1983), no.~2, 687--708.

\bibitem[GL87]{GordonLuecke}
C.~McA. Gordon and J. Luecke,
\emph{Only integral Dehn surgeries can yield reducible manifolds},
Math. Proc. Cambridge Philos. Soc. \textbf{102} (1987), no.~1, 97--101.

\bibitem[HLZ15]{HLZ}
J. Hom, T. Lidman, and N. Zufelt,
\emph{Reducible surgeries and Heegaard Floer homology},
Math. Res. Lett. \textbf{22} (2015), no.~3, 763--788.

\bibitem[How02]{Howie}
J. Howie,
\emph{A proof of the Scott--Wiegold conjecture on free products of
cyclic groups},
J. Pure Appl. Algebra \textbf{173} (2002), no.~2, 167--176.

\bibitem[Mil62]{Milnor}
J. Milnor,
\emph{A unique decomposition theorem for \(3\)-manifolds},
Amer. J. Math. \textbf{84} (1962), no.~1, 1--7.

\bibitem[MT07]{MorganTian}
J. Morgan and G. Tian,
\emph{Ricci Flow and the Poincar\'e Conjecture},
Clay Mathematics Monographs, vol.~3,
American Mathematical Society, Providence, RI, and
Clay Mathematics Institute, Cambridge, MA, 2007.

\bibitem[Mos71]{Moser}
L. Moser,
\emph{Elementary surgery along a torus knot},
Pacific J. Math. \textbf{38} (1971), no.~3, 737--745.

\bibitem[Sch90]{Scharlemann}
M. Scharlemann,
\emph{Producing reducible \(3\)-manifolds by surgery on a knot},
Topology \textbf{29} (1990), no.~4, 481--500.

\bibitem[Sta65]{Stallings}
J.~R. Stallings,
\emph{A topological proof of Grushko's theorem on free products},
Math. Z. \textbf{90} (1965), 1--8.

\end{thebibliography}

\end{document}
